\section{Introduction}

Large Language Models (LLMs) have become the backbone of modern AI applications~\cite{brown2020language}.
However, their deployment at scale remains expensive due to massive computational and memory requirements.
A key bottleneck in LLM inference is the key-value (KV) cache, which stores attention keys and values for all previously generated tokens to avoid recomputation~\cite{pope2023efficiently}.
In long-context and large-batch scenarios, the KV cache can exceed model parameters by multiple times: for 540B PaLM with batch size 512 and context length 2048, KV cache alone consumes 3TB of memory~\cite{pope2023efficiently}.

To address this, quantization techniques reduce the bit-width of KV cache entries.
\kivi{}~\cite{kivi} proposes tuning-free asymmetric 2-bit quantization, achieving $8\times$ memory savings with per-channel quantization for keys and per-token quantization for values.
However, while \kivi{} effectively reduces KV cache \emph{size}, the dequantization overhead during attention computation introduces latency that partially offsets memory bandwidth savings.

The root cause is the mathematical expansion of asymmetric quantized dot products:
\begin{equation}
Q \cdot K_i \approx s_Q s_K \left[\sum_{j=1}^{d} \hat{q}_j \hat{k}_{ij} - z_K\sum_{j=1}^{d} \hat{q}_j - z_Q\sum_{j=1}^{d} \hat{k}_{ij} + d z_Q z_K\right]
\label{eq:asym_expand}
\end{equation}
This reveals four distinct operations beyond the standard dot product $\sum \hat{q}\hat{k}$: two channel-wise sums ($\sum\hat{q}$, $\sum\hat{k}$) and a constant correction.
On GPU, these extra operations add significant overhead.

Processing-in-memory (PIM) offers a promising alternative by placing compute units close to memory~\cite{mutlu2020processing,seshadri2017ambit}.
Prior PIM designs for attention~\cite{duplex} used logic-die PIM, which places compute on a separate die connected via TSVs, introducing bandwidth bottlenecks.
Near-bank PIM, where lightweight PEs reside inside each DRAM bank, provides higher local bandwidth and lower latency for bank-local operations.

In this paper, we propose \pimattn{}, a heterogeneous attention framework that decomposes the asymmetric quantized attention across near-bank PIM and GPU.
Our key insight: the operations in Eq.~\ref{eq:asym_expand} have fundamentally different computational characteristics.
The main dot product $\sum\hat{q}\hat{k}$ and reduction sums operate on 2-bit data --- ideal for near-bank PIM with efficient 2-bit PEs.
The score-value ($S\cdot V$) operation involves float16 tensors --- better suited for GPU's massive float16 throughput.
Softmax is lightweight and stays on GPU.

\pimattn{} maps the decomposed Q@K operations to near-bank PIM with 2-bit compute units and routes score@V to GPU.
We implement a cycle-accurate simulator with per-stage timing breakdown and configurable PE microarchitecture.
Evaluation on Llama-3-8B with 16K context demonstrates $2.1\mu s$ per-sequence Q@K latency on 2-bit KV with 2048 PIM banks --- $1.4\times$ faster than GPU.
We further show that PIM-side dequantization incurs $2\times$ overhead, validating direct 2-bit compute.

Our contributions:
\begin{itemize}
\item We identify the computational mismatch between asymmetric quantized attention and existing hardware, motivating a heterogeneous decomposition.
\item We design \pimattn{}, mapping decomposed Q@K to near-bank PIM with 2-bit PEs, while routing float16-heavy operations to GPU.
\item We build a cycle-accurate PIM simulator with configurable architecture and per-stage timing verification (ratio 1.00$\times$).
\item We provide comprehensive experimental analysis across six configurations.
\end{itemize}
